\chapter{Metodologia}
\label{chap:metodologia}

Este capítulo descreve os métodos experimentais empregados na caracterização do sistema de
propulsão do automodelo. A apresentação segue a cadeia da medida: o princípio do ensaio de
tração e a sua montagem (Seções \ref{sec:montagem} e \ref{sec:angulocabo}), a determinação do
raio efetivo da roda (Seção \ref{sec:raioefetivo}), a instrumentação elétrica e o registro dos
dados (Seção \ref{sec:cadeia}), o procedimento experimental (Seção \ref{sec:procedimento}) e,
por fim, as incertezas de medição (Seção \ref{sec:incertezas}). Para cada escolha, apresentam-se
o critério que a fundamenta e as alternativas consideradas, de modo a permitir a reprodução
integral do ensaio.

\section{O Ensaio de Tração Estática}
\label{sec:montagem}

Caracterizar o sistema de propulsão exige medir, simultaneamente, a força de tração
desenvolvida pelas rodas e a corrente elétrica consumida pelo conjunto, uma vez que a curva que
este trabalho busca estabelecer relaciona o torque na roda à corrente de alimentação.

A tração é a força com que o veículo é impulsionado sobre o solo --- o efeito mecânico final
de toda a cadeia motor--transmissão--roda. Diretamente, ela só se manifesta enquanto o veículo
se desloca; em condição estacionária, contudo, a mesma força surge de outro modo. Se o
veículo for impedido de se mover, o pneu não consegue rolar e a tração que seria entregue ao
solo transfere-se integralmente para a estrutura que retém o veículo. O ensaio de tração
estática explora exatamente essa condição: o automodelo é ancorado, e a reação da ancoragem ---
medida como a tensão no cabo que o prende --- iguala a tração desenvolvida pelas rodas,
dispensando um dinamômetro dedicado. Como o veículo está parado, a tração medida corresponde
à condição de partida, em que o motor opera em baixa rotação e máxima exigência de torque ---
precisamente a região de maior interesse para o controle de aceleração.

A leitura da balança empregada no ensaio é expressa em quilogramas, mas o que o instrumento
mede quando o seu gancho é tracionado é uma força. O valor exibido é a massa cuja força-peso
equilibraria a tensão aplicada ao gancho --- o quilograma-força, definido como o peso de um
quilograma em condições normais de gravidade e equivalente a $9{,}81$ N. Como o gancho é puxado
horizontalmente pela reação da ancoragem, e não sustenta massa alguma na vertical, a leitura em
quilogramas representa força de tração, e não peso. Ao longo do trabalho, as leituras são
convertidas para newtons por meio de $F\,[\mathrm{N}] = 9{,}81\,m\,[\mathrm{kgf}]$. O torque na
roda é obtido da força pela roda efetiva, determinada na Seção \ref{sec:raioefetivo}, e a
corrente, pela cadeia de aquisição descrita na Seção \ref{sec:cadeia}.

\section{Geometria do Cabo e o Efeito Cosseno}
\label{sec:angulocabo}

A força que equilibra a tração do veículo atua na direção horizontal --- a direção em que o
veículo tentaria se deslocar. A tensão que a balança registra, porém, atua ao longo do cabo, e
o cabo não é perfeitamente horizontal: o engate no veículo se encontra a $28$ cm do solo,
ao passo que a balança se encontra a $15$ cm. O cabo e a linha horizontal entre os mesmos dois
pontos formam, com a diferença de altura, um triângulo retângulo, cuja hipotenusa é o próprio
cabo e cujo cateto vertical, oposto ao ângulo $\theta$ entre o cabo e a horizontal, é a
diferença de altura entre os engates.

Cada uma das três dimensões da montagem foi medida três vezes, e os valores adotados ---
comprimento do cabo de $126$ cm, alturas de $15$ cm e $28$ cm --- são as médias das três
leituras. A diferença de altura, de $28 - 15 = 13$ cm, sobre o comprimento do cabo, define,
pela razão entre o cateto oposto e a hipotenusa,
\begin{equation}\label{eq:sen_theta}
\mathrm{sen}\,\theta = \frac{13}{126}
\qquad\Rightarrow\qquad
\theta = 5{,}9^\circ.
\end{equation}
A componente horizontal da tensão --- a que efetivamente equilibra a tração --- é $F\cos\theta$;
tomar a leitura da balança diretamente como força de tração equivale a aproximar $\cos\theta$
por $1$, com erro relativo dado pela fração da tensão que foge à direção horizontal,
\begin{equation}\label{eq:erro_cosseno}
1 - \cos\theta = 1 - 0{,}995 = 0{,}0053 = 0{,}53\%.
\end{equation}
O desvio sistemático máximo decorrente do ângulo é, portanto, de $0{,}53\%$; em vez de corrigi-lo
em cada leitura, esse valor é carregado como componente sistemática no orçamento de incertezas
da Seção \ref{sec:incertezas}.

\section{Medida do Raio Efetivo da Roda}
\label{sec:raioefetivo}

O torque na roda é obtido da força de tração por meio do raio em que essa força atua --- a
distância entre o centro da roda e o contato do pneu com o solo. Esse raio, contudo, não é o
nominal do pneu novo e descarregado, mas o raio efetivo de rolamento, medido sob as condições
de carga do veículo: o pneu se deforma sob o peso do conjunto, e o raio com que ele efetivamente
trabalha é menor que o nominal.

O raio efetivo foi medido pelo método de \emph{roll-out}. Marcaram-se o pneu e o solo em um
mesmo alinhamento; o veículo, encostado em uma parede para garantir que o deslocamento se
desse em linha reta, foi movimentado até completar cinco voltas da roda. A distância percorrida
entre as marcas do solo corresponde, nessas condições, a cinco circunferências de rolamento.
O procedimento foi repetido três vezes, e a distância adotada, $L = 236$ cm, é a média das
três leituras. O raio efetivo segue diretamente da circunferência:
\begin{equation}\label{eq:raio_efetivo}
r_\mathrm{ef} = \frac{L}{2\pi n}
              = \frac{236}{2\pi \times 5}
              = 7{,}51\ \mathrm{cm}
              = 0{,}0751\ \mathrm{m},
\end{equation}
em que $n = 5$ é o número de voltas completas.

O valor nominal do fabricante para o pneu novo, sem carga, é de $80$ mm de raio; o raio
efetivo medido, de $75{,}1$ mm, é $6{,}1\%$ menor --- a deflexão do pneu sob a carga do
conjunto. Fecha-se, assim, a cadeia anunciada na Seção \ref{sec:montagem}: o torque na roda é
\begin{equation}\label{eq:torque_f_r}
T_\mathrm{roda} = F\,r_\mathrm{ef},
\end{equation}
e o emprego do raio nominal em lugar do efetivo injetaria erro sistemático de $6{,}1\%$ em
todos os pontos da curva característica.

\clearpage
